\documentclass[10pt,a4paper]{amsart}
\usepackage[margin=19mm]{geometry}
\usepackage{amsmath,amssymb,mathtools}
\usepackage[scr=rsfs,cal=euler]{mathalfa}
\usepackage{xfrac}
\usepackage[unicode,hidelinks,bookmarks=false]{hyperref}
\newcommand{\ie}{i.e.\ }
\newcommand{\cf}{cf.\ }
\newcommand{\resp}{respectively\ }
\newcommand{\Subsection}{\subsection}
\providecommand{\nobreakdash}{\nobreak\hbox{-}}
\setlength{\emergencystretch}{2em}
\multlinegap0pt
\usepackage{amssymb}
\newtheorem{proposition}{Proposition}
\title{Properties of two Chebyshev-like polynomial sequences}
\author{Karl Dilcher, Seon-Hong Kim,, Kenneth B. Stolarsky}
\date{Source: arXiv:2607.16940v1 (CC BY 4.0)}
\begin{document}
\maketitle

\noindent\emph{Source and licence.} Properties of two Chebyshev-like polynomial sequences. Version arXiv:2607.16940v1 (CC BY 4.0), distributed by arXiv under the Creative Commons Attribution 4.0 International licence, http://creativecommons.org/licenses/by/4.0/, as recorded in the arXiv licence metadata for this exact version and shown on the abstract page https://arxiv.org/abs/2607.16940v1. Full licence notice: \texttt{licenses/arxiv-2607.16940-CC-BY-4.0.txt}.\par\smallskip

\noindent\textit{Editorial scope.} Counted unit: the article's proposition prop:5.4 (source0.tex lines 833--838). The article's complete original proof of it is delivered with it (source0.tex lines 855--928, one \texttt{proof} environment of 74 lines, which the article places away from the statement). No mathematics is written, repaired or re-derived: the statement and the proof are the article's own lines, in the article's own notation, and no retained line is substituted or re-flowed.  Retained uncounted as the source-local context the proof needs: lines 182--184; lines 244--246; lines 824--827; lines 844--847.  Nothing was omitted from any retained span, so the package carries no omission record.  No retained line contains a citation, so the wrapper carries no bibliography and no bibliography entry.  No retained line contains a figure environment, an image-inclusion command or drawing code, so no image is delivered and the package declares no asset dependency.  Editorial formatting only, in addition: an AMS \texttt{amsart} preamble that restates the article's own declarations and macros used by the retained lines, namely the environments \texttt{proposition} on separate counters, together with the standard packages those lines need.  The retained environments print in this document as Proposition 1 in order of appearance, whereas the article prints the same items with the numbering of its own full document.  The complete native source of the article as distributed in the arXiv e-print for this exact version is bundled unchanged as \texttt{sources/205518/source0.tex}.
\begin{equation}\label{2.6}
p_n(1) = (n+1)2^{2n},\qquad p_n(-1) = 2^{2n}.
\end{equation}

\begin{equation}\label{2.10}
p_n(z)=\frac{(z-1)^{2n+4}-\left(z^2-(4n+6)z+1\right)(z+1)^{2n+2}}{16z^2}.
\end{equation}

\begin{equation}\label{5.18}
{\rm Disc}(f)=(-1)^{\frac{\mu(\mu-1)}{2}}a_{\mu}^{\mu-2}
\prod_{j=1}^{\mu}f'(\alpha_j),
\end{equation}

\begin{equation}\label{5.20}
p_n'(z) = \frac{2n}{z-1}\cdot p_n(z)-4(z-1)^{2n-2}
\sum_{k=1}^n k(k+1)\left(\frac{z+1}{z-1}\right)^{2k-1}.
\end{equation}

\begin{proposition}\label{prop:5.4}
For all $n\geq 1$ we have
\begin{equation}\label{5.19}
{\rm Disc}(p_n) = (-1)^n 2^{4n^2+2}(n+1)^{2n-3}(n+2)^{2n-2}.
\end{equation}
\end{proposition}

\begin{proof}[Proof of Proposition~\ref{prop:5.4}]
We use \eqref{5.18} with $f=p_n, \mu=2n$, and $a_{\mu}=(n+1)(n+2)/2$. Then
\begin{equation}\label{5.21}
{\rm Disc}(p_n)=(-1)^n \left(\frac{(n+2)(n+1)}{2}\right)^{2n-2}
\cdot\prod_{i=1}^{2n}p_n'(\alpha_i),
\end{equation}
where $\alpha_1, \ldots, \alpha_{2n}$ are the zeros of $p_n(z)$. Now, by
\eqref{5.20} we have
\begin{equation}\label{5.22}
p_n'(\alpha_i) =-4\cdot \frac{(\alpha_i-1)^{2n-1}}{\alpha_i+1}\cdot 
\sum_{k=1}^n k(k+1)\left(\frac{\alpha_i+1}{\alpha_i-1}\right)^{2k}.
\end{equation}
To evaluate the sum in \eqref{5.22}, we begin with the identity
\[
1+y+y^2+\cdots+y^{n+1} = \frac{1-y^{n+2}}{1-y}.
\]
After differentiating twice, we get
\begin{equation}\label{5.23}
\sum_{k=1}^n k(k+1)y^{k-1} = \frac{1}{(1-y)^2}\left(2\cdot\frac{1-y^{n+1}}{1-y}
-(n+1)(n+2)y^n+n(n+1)y^{n+1}\right).
\end{equation}
Next, we set $z=\alpha_i$ in \eqref{2.10} for any $i=1,2, \ldots, 2n$. 
Writing $\alpha=\alpha_i$ for convenience of notation, we get
\begin{equation}\label{5.23a}
(\alpha-1)^2
=(\alpha^2-(4n+6)\alpha+1)\left(\frac{\alpha+1}{\alpha-1}\right)^{2(n+1)},
\end{equation}
which can be rewritten as
\[
(\alpha-1)^2-(\alpha-1)^2\left(\frac{\alpha+1}{\alpha-1}\right)^{2(n+1)}
=(n+1)\left((\alpha-1)^2-(\alpha+1)^2\right)
\left(\frac{\alpha+1}{\alpha-1}\right)^{2(n+1)},
\]
and after some further straightforward manipulation we get
\begin{equation}\label{5.24}
\frac{1-y^{n+1}}{1-y}=(n+1)y^{n+1},\qquad\hbox{where}\quad
y= \left(\frac{\alpha+1}{\alpha-1}\right)^2.
\end{equation}
Substituting this into the right-hand side of \eqref{5.23}, we get
\begin{align}
\sum_{k=1}^n k(k+1)y^{k-1} &= \frac{1}{(1-y)^2}\left(2(n+1)y^{n+1}
-(n+1)(n+2)y^n+n(n+1)y^{n+1}\right) \label{5.25}\\
&=-\frac{(n+1)(n+2)}{1-y}y^{n},\nonumber
\end{align}
where $y$ is as given in \eqref{5.24}. By \eqref{5.23a}, an easy calculation 
gives
\begin{equation}\label{5.26}
\frac{1}{1-\left(\frac{\alpha+1}{\alpha-1}\right)^2}\cdot
\left(\frac{\alpha+1}{\alpha-1}\right)^{2n+2}
=\frac{(\alpha+1)^{2n+2}}{-4\alpha(\alpha-1)^{2n}}.
\end{equation}
With \eqref{5.22}, \eqref{5.25} and \eqref{5.26} we then get
\[
p_n'(\alpha_i) 
=-(n+1)(n+2)\cdot\frac{(\alpha_i+1)^{2n+1}}{\alpha_i(\alpha_i-1)},
\]
and with \eqref{5.21},
\begin{equation}\label{5.28}
{\rm Disc}(p_n)=(-1)^n \frac{((n+1)(n+2))^{4n-2}}{2^{2n-2}}
\cdot\prod_{i=1}^{2n}\frac{(\alpha_i+1)^{2n+1}}{\alpha_i(\alpha_i-1)}.
\end{equation}
To evaluate the product in \eqref{5.28}, we use the fact that the identity
\[
p_n(z)=\frac{(n+1)(n+2)}{2}\cdot\prod_{i=1}^{2n}(z-\alpha_i),
\]
together with \eqref{2.6}, implies
\[
\prod_{i=1}^{2n}(1-\alpha_i)=\frac{2^{2n+1}}{n+2},\qquad
\prod_{i=1}^{2n}(1+\alpha_i)=\frac{2^{2n+1}}{(n+1)(n+2)},\qquad
\prod_{i=1}^{2n}\alpha_i=1.
\]
Finally, these identities, together with \eqref{5.28}, give \eqref{5.19} after
collecting the powers of 2, $n+1$ and $n+2$.
\end{proof}

\end{document}
